\section{Alternating harmonic sums: complete even-weight evaluation}
\label{sec:euler}

The Gaussian depth draft studies
\begin{equation}
S_p=\sum_{n=0}^{\infty}\frac{(-1)^nH_n}{(2n+1)^p},
\qquad H_0=0,\qquad H_n=\sum_{k=1}^{n}\frac1k.
\label{eq:euler-definition}
\end{equation}
For ordinary convergence the correct index range is \(p\ge1\).
At \(p=0\), the summands do not tend to zero. An Abel value at that
index would be a separate definition, not the sum in
\eqref{eq:euler-definition}.

The odd-index evaluations below are a specialization of the existing
Euler-sum literature. We prove them independently because the integral
argument explains exactly which symmetry supplies the evaluation and
leads to a useful global generating function. Neither the evaluation
nor the symmetry supplies a nonmembership theorem for the even indices.

\subsection{Convergence and the Mellin representation}

\begin{lemma}\label{lem:euler-mellin}
For every integer \(p\ge1\), the series \eqref{eq:euler-definition}
converges and satisfies \(|S_p|\le3^{-p}\). It converges absolutely
precisely when \(p>1\). With
\[
Q(x)=\frac{\log(1+x^2)}{1+x^2},
\]
one has
\begin{equation}
S_p=-\frac1{(p-1)!}\int_0^1(-\log x)^{p-1}Q(x)\,dx.
\label{eq:euler-mellin}
\end{equation}
\end{lemma}

\begin{proof}
For \(n\ge1\),
\[
\frac{H_{n+1}}{H_n}
=1+\frac1{(n+1)H_n}
\le1+\frac1{n+1}
<\frac{2n+3}{2n+1}.
\]
Thus \(H_n/(2n+1)^p\) decreases to zero for every \(p\ge1\).
The alternating-series theorem gives convergence and the bound by
the first nonzero term, \(3^{-p}\). The estimate
\(H_n\le1+\log n\) proves absolute convergence for \(p>1\);
for \(p=1\), even the lower bound \(H_n\ge1\) proves divergence
of the absolute series.

To justify the endpoint \(p=1\) together with the other indices, take
\(0\le r<1\). The harmonic-number generating identity gives
\[
\sum_{n\ge0}(-r)^nH_nx^{2n}
=-\frac{\log(1+rx^2)}{1+rx^2},\qquad 0\le x\le1.
\]
The series is uniformly absolutely convergent in \(x\), and
\[
\int_0^1x^{2n}(-\log x)^{p-1}\,dx
=\frac{(p-1)!}{(2n+1)^p}.
\]
Termwise integration therefore proves the regularized version of
\eqref{eq:euler-mellin}. Since
\[
0\le \frac{\log(1+rx^2)}{1+rx^2}\le x^2,
\]
dominated convergence gives the stated integral as \(r\uparrow1\).
Abel's theorem identifies the series limit with the already established
ordinary sum. In particular, no unproved interchange of a conditionally
convergent boundary series is required.
\end{proof}

\subsection{Evaluation at every odd index}

\begin{theorem}[All odd-index harmonic sums]\label{thm:euler-odd}
For every integer \(m\ge0\),
\begin{equation}
\boxed{\displaystyle
S_{2m+1}=(2m+1)\beta(2m+2)-2\beta(2m+1)\log2
-\sum_{j=1}^{m}(2-2^{-2j})
 \beta(2m-2j+1)\zeta(2j+1).}
\label{eq:euler-odd}
\end{equation}
The sum is empty for \(m=0\).
\end{theorem}

\begin{proof}
We first relate a half-line logarithmic moment to a full-line moment.
Put
\[
J_m=\int_0^1\log^{2m}x\,Q(x)\,dx,
\qquad
L_m=\int_0^{\infty}\log^{2m}x\,Q(x)\,dx.
\]
All these integrals converge: \(Q(x)=O(x^2)\) at zero and
\(Q(x)=O(x^{-2}\log x)\) at infinity.
Substitution \(x=1/u\) in the part over \((1,\infty)\) uses the
exact identity
\[
\frac{Q(1/u)}{u^2}
=Q(u)-\frac{2\log u}{1+u^2}.
\]
Consequently,
\begin{align}
L_m
&=2J_m-2\int_0^1\frac{\log^{2m+1}u}{1+u^2}\,du\notag\\
&=-2(2m)!S_{2m+1}+2(2m+1)!\beta(2m+2).
\label{eq:euler-inversion}
\end{align}
Here the first term follows from Lemma~\ref{lem:euler-mellin}; the
second follows by expanding \((1+u^2)^{-1}\). That integration is
absolutely summable, since the integrated terms have size
\((2m+1)!/(2n+1)^{2m+2}\).

It remains to evaluate \(L_m\). Define the full Mellin transform
\[
F(a)=\int_0^{\infty}x^{a-1}Q(x)\,dx.
\]
The integral in fact converges for \(-2<\Repart a<2\). On the
smaller strip \(0<\Repart a<2\), the beta integral gives
\begin{equation}
F(a)=\frac{\pi}{2\sin(\pi a/2)}
       \bigl[\psi(1)-\psi(1-a/2)\bigr],
\label{eq:euler-full-mellin}
\end{equation}
where \(\psi=\Gamma'/\Gamma\). Indeed, for
\(\Repart b>\Repart a/2>0\),
\[
\int_0^{\infty}x^{a-1}(1+x^2)^{-b}\,dx
=\frac{\Gamma(a/2)\Gamma(b-a/2)}{2\Gamma(b)}.
\]
Differentiate with respect to \(b\), negate the result, and set
\(b=1\). Differentiation under the integral is valid uniformly on
compact subsets of this parameter region, because the additional
factor \(\log(1+x^2)\) preserves the integrable bounds at zero
and infinity. Euler's reflection formula then yields
\eqref{eq:euler-full-mellin}. The same bounds with further powers
of \(\log x\) justify differentiation in \(a\), so
\begin{equation}
F^{(2m)}(1)=L_m.
\label{eq:euler-moments}
\end{equation}

The digamma duplication formula and the convergent polygamma series
give
\[
\psi(1)-\psi(1/2)=2\log2,
\qquad
\psi^{(k)}(1/2)
=(-1)^{k+1}k!(2^{k+1}-1)\zeta(k+1),\quad k\ge1.
\]
For clarity, the second identity follows by separating the odd terms
in \(\psi^{(k)}(z)=(-1)^{k+1}k!\sum_{n\ge0}(n+z)^{-k-1}\).
Expanding \eqref{eq:euler-full-mellin} at \(a=1\), its even part
is therefore
\begin{equation}
\frac{F(1+t)+F(1-t)}2
=\pi\sec\!\left(\frac{\pi t}{2}\right)
 \left[\log2+\sum_{j\ge1}
 (1-2^{-2j-1})\zeta(2j+1)t^{2j}\right],
\qquad |t|<1.
\label{eq:euler-even-mellin}
\end{equation}

The needed coefficients of the secant can be obtained without any
assumption about special values. Another beta integral and the same
inversion give
\[
\frac{\pi}{2}\sec\!\left(\frac{\pi t}{2}\right)
=\int_0^1\frac{x^t+x^{-t}}{1+x^2}\,dx
=2\sum_{r\ge0}\beta(2r+1)t^{2r},\qquad |t|<1.
\]
To justify the last equality, expand the numerator in powers of \(t\)
under an integrable bound \(x^{-|t|}\). The logarithmic moments
with \(r\ge1\) may be integrated termwise absolutely; for \(r=0\)
the identity is \(\arctan1=\beta(1)=\pi/4\).
Taking the coefficient of \(t^{2m}\) in
\eqref{eq:euler-even-mellin} now yields
\[
\frac{L_m}{(2m)!}
=4\beta(2m+1)\log2
 +4\sum_{j=1}^{m}(1-2^{-2j-1})
       \beta(2m-2j+1)\zeta(2j+1).
\]
Substituting this into \eqref{eq:euler-inversion} and dividing by
\(2(2m)!\) proves \eqref{eq:euler-odd}.
\end{proof}

\paragraph{Literature identification and normalization.}
In Xu and Wang's notation for linear Euler \(R\)-sums,
\[
R_{1,q}^{1,-1}
=\sum_{n\ge1}\frac{(-1)^{n-1}H_{n-1}}{(n-1/2)^q}
=2^qS_q.
\]
For odd \(q\ge3\), equation~(20) on printed page~994 of
\cite{XuWang} specializes to \eqref{eq:euler-odd}.
The conversion uses their conventions
\[
\widetilde t(k;-1)=-2^k\beta(k),\qquad
\widetilde t(k;1)=(2^k-1)\zeta(k)\quad(k\ge2),
\qquad \widetilde t(1;1)=2\log2.
\]
The \(m=0\) case follows directly from the proof above. Thus the
theorem is an explicit, independently proved specialization of known
results, not a claim of a globally new evaluation theorem.

For example,
\begin{align*}
S_1&=G-\frac\pi2\log2,\\
S_3&=3\beta(4)-\frac{\pi^3}{16}\log2
                      -\frac{7\pi}{16}\zeta(3),\\
S_9&=9\beta(10)-\frac{1385\pi^9}{20643840}\log2
       -\frac{427\pi^7}{737280}\zeta(3)
       -\frac{155\pi^5}{24576}\zeta(5)\\
   &\hspace{1.7em}-\frac{127\pi^3}{2048}\zeta(7)
       -\frac{511\pi}{1024}\zeta(9).
\end{align*}
Every odd beta value is a rational multiple of the corresponding odd
power of \(\pi\), as also follows by comparing the secant coefficients
above. Hence \eqref{eq:euler-odd} is homogeneous of weight \(2m+2\)
in the single-value generators used in the draft. In particular, it
proves the positive, even-weight half of the proposed alternation
pattern at all weights.

\subsection{A meromorphic generating function and its exact singularities}

\begin{theorem}[Even generator]\label{thm:euler-generator}
The power series
\[
E(t)=\sum_{m\ge0}S_{2m+1}t^{2m}
\]
has radius of convergence exactly \(3\). It extends to an even
meromorphic function on \(\C\), with the closed form
\begin{equation}
\boxed{\displaystyle
E(t)=\frac{\psi_1((1+t)/4)-\psi_1((3+t)/4)}{16}
 +\frac{\pi\,[\psi((1+t)/2)-\psi(1)]}
        {4\cos(\pi t/2)},}
\label{eq:euler-generator}
\end{equation}
where \(\psi_1=\psi'\), and apparent singularities are interpreted
by meromorphic continuation. Its only poles are simple and occur at
\(t=\pm(2n+1)\), \(n\ge1\), with
\begin{equation}
\Res_{t=2n+1}E(t)=\frac{(-1)^{n+1}H_n}{2},\qquad
\Res_{t=-(2n+1)}E(t)=\frac{(-1)^nH_n}{2}.
\label{eq:euler-residues}
\end{equation}
In particular, the apparent singularities at \(t=\pm1\) are
removable, and
\begin{equation}
E(1)=E(-1)=-\frac{\pi^2}{48}.
\label{eq:euler-removable}
\end{equation}
\end{theorem}

\begin{proof}
Lemma~\ref{lem:euler-mellin} and the bound \(|S_p|\le3^{-p}\)
first define \(E\) on \(|t|<3\). Summing the even exponential
series under the Mellin integral gives
\begin{equation}
E(t)=-\frac12\int_0^1(x^t+x^{-t})Q(x)\,dx.
\label{eq:euler-generator-integral}
\end{equation}
For \(|t|<3\), absolute values of the summed integrand are bounded
by \(Q(x)\cosh(|t|\log x)\), which is integrable at zero since
\(Q(x)=O(x^2)\). The right-hand side of
\eqref{eq:euler-generator-integral} is holomorphic on the larger
vertical strip \(|\Repart t|<3\): on each compact subset, the same
argument applies with a fixed exponent strictly smaller than \(3\),
and all parameter derivatives add only powers of \(|\log x|\).
It also proves evenness there.

To evaluate the integral, let
\[
J(t)=\int_0^1x^tQ(x)\,dx,
\qquad
K(t)=\int_0^1\frac{x^t\log x}{1+x^2}\,dx.
\]
For \(|\Repart t|<1\), splitting the full Mellin integral at one
and inverting the second part gives
\[
F(1-t)=J(-t)+J(t)-2K(t).
\]
In this strip, expansion of \((1+x^2)^{-1}\) and absolute
summability after integration give
\begin{align*}
K(t)
&=-\sum_{n\ge0}\frac{(-1)^n}{(2n+1+t)^2}\\
&=-\frac1{16}
 \bigl[\psi_1((1+t)/4)-\psi_1((3+t)/4)\bigr].
\end{align*}
Since \(E(t)=-F(1-t)/2-K(t)\), substitution of
\eqref{eq:euler-full-mellin} proves
\eqref{eq:euler-generator} in \(|\Repart t|<1\).

The exact global singularity statement should not be inferred merely
from the individual terms of that formula, because they have cancelling
poles. Instead, retain the conditionally convergent coefficient
\(S_1\) separately and sum all the absolutely convergent higher
coefficients. For \(|t|<3\), this gives
\begin{equation}
E(t)=S_1+\sum_{n\ge1}
 \frac{(-1)^nH_n t^2}
 {(2n+1)((2n+1)^2-t^2)}.
\label{eq:euler-generator-partial-fractions}
\end{equation}
The interchange is valid because
\[
\sum_{n\ge1}\frac{H_n}{2n+1}
 \sum_{m\ge1}\left(\frac{|t|}{2n+1}\right)^{2m}
=\sum_{n\ge1}
 \frac{H_n|t|^2}{(2n+1)((2n+1)^2-|t|^2)}<\infty.
\]
For any compact set avoiding the displayed poles in
\eqref{eq:euler-generator-partial-fractions}, the summands for
sufficiently large \(n\) are uniformly
\(O((1+\log n)/n^3)\). Thus the series converges normally there
and defines a meromorphic function on \(\C\). At a point
\(2n+1\) or \(-(2n+1)\), precisely one summand is singular;
its elementary residue is \eqref{eq:euler-residues}. These
residues are nonzero, and there are no further poles. The right-hand
side of \eqref{eq:euler-generator} is also meromorphic on \(\C\)
and agrees on an open strip, so the meromorphic identity theorem
proves that the two representations agree everywhere. In particular,
all additional apparent poles in the polygamma expression cancel.
The nearest actual poles to zero are \(3\) and \(-3\), proving
that the Taylor radius is exactly \(3\).

Finally, \eqref{eq:euler-generator-integral} can be evaluated directly
at \(t=1\), inside its holomorphy strip:
\[
E(1)=-\frac12\int_0^1\frac{\log(1+x^2)}{x}\,dx
     =\frac14\Li_2(-1)=-\frac{\pi^2}{48}.
\]
Evenness gives the same value at \(-1\), proving
\eqref{eq:euler-removable} without a delicate cancellation of
separate Laurent expansions.
\end{proof}

\subsection{The full generator and the unresolved parity direction}

The symmetry can be stated more precisely by keeping all indices.

\begin{proposition}\label{prop:euler-full-generator}
The full generating function
\[
A(t)=\sum_{p\ge1}S_pt^{p-1}
\]
has Taylor radius \(3\), integral representation
\begin{equation}
A(t)=-\int_0^1x^{-t}Q(x)\,dx,\qquad \Repart t<3,
\label{eq:euler-full-generator-integral}
\end{equation}
and global meromorphic representation
\begin{equation}
A(t)=S_1+\sum_{n\ge1}
 \frac{(-1)^nH_n t}{(2n+1)(2n+1-t)}.
\label{eq:euler-full-generator-poles}
\end{equation}
Its only poles are simple, at \(t=2n+1\), \(n\ge1\), with
residue \((-1)^{n+1}H_n\). Moreover,
\begin{equation}
\frac{A(t)+A(-t)}2=E(t),\qquad
\frac{A(t)-A(-t)}2=\sum_{m\ge1}S_{2m}t^{2m-1}
\quad (|t|<3).
\label{eq:euler-parity-generators}
\end{equation}
\end{proposition}

\begin{proof}
For \(|t|<3\), insert \eqref{eq:euler-mellin} and sum the
exponential series \(\sum_{p\ge1}t^{p-1}(-\log x)^{p-1}/(p-1)!
=x^{-t}\). The absolute sum is bounded by \(x^{-|t|}Q(x)\),
which is integrable. The resulting integral is holomorphic for
\(\Repart t<3\), by the same compact-set bounds as above.
Separating \(p=1\) and summing the geometric series for \(p\ge2\)
gives \eqref{eq:euler-full-generator-poles}; the needed absolute
bound is
\[
\sum_{n\ge1}\frac{H_n|t|}{(2n+1)(2n+1-|t|)}<\infty,
\qquad |t|<3.
\]
On a compact set avoiding positive odd integers at least three,
the resulting rational summands are uniformly
\(O((1+\log n)/n^2)\). Normal convergence gives the stated global
continuation. The single summand singular at \(t=2n+1\) has residue
\((-1)^{n+1}H_n\), and that residue is nonzero. This proves both
the complete pole list and the exact radius. Taking even and odd
parts of the Taylor series proves \eqref{eq:euler-parity-generators}.
\end{proof}

The inversion \(x\mapsto1/x\) determines the even part of \(A\)
through the full Mellin transform. It leaves the odd part as a
different analytic quantity. The absence of the same reduction for
that part is not a proof that its coefficients lie outside the algebra
generated by single polylogarithmic values. In particular, no result
here proves nonmembership for \(S_{2m}\), which have odd total
weight. Such a claim requires an additional arithmetic or motivic
argument for the specified constants and specified algebra. The
presence of an even beta value likewise supplies no independence
certificate: even the irrationality of \(G=\beta(2)\) is not known.
The correction to the draft is therefore precise: replace the proposed
two-sided alternation theorem by \eqref{eq:euler-odd} and its proved
even-weight membership statement, and retain the opposite direction
as a research question.
